\usepackage{multirow}
\usepackage{tabularx}
\renewcommand\tabularxcolumn[1]{m{#1}} % Centrado vertical en tabularx
\usepackage{array}

\geometry{a4paper, margin=1in}

\setlength{\parindent}{0pt}
\setlength{\parskip}{0pt}
\renewcommand{\arraystretch}{1.3}


\chapter{Control de Acceso a Actividades}

%======================================== Caso de uso 29 ====================
\begin{table}[h!]
\centering
\begin{tabularx}{\textwidth}{|l|X|X|}
\hline
\textbf{Caso de Uso} & Registrar acceso a actividad & CU-29 \\ \hline
\textbf{Actores} & \multicolumn{2}{l|}{Responsable (I), Cliente} \\ \hline
\textbf{Tipo} & \multicolumn{2}{l|}{Primario, Esencial} \\ \hline
\textbf{Referencias} & \multicolumn{2}{l|}{} \\ \hline
\textbf{Precondición} & \multicolumn{2}{l|}{La actividad debe de existir en el sistema.} \\ \hline
\textbf{Poscondición} & \multicolumn{2}{l|}{Queda registrado el acceso del cliente en el sistema.} \\ \hline
\textbf{Autor} & \multicolumn{2}{l|}{Julian Carrion Tovar \hfill Fecha \hfill Versión 1.0} \\ \hline
\end{tabularx}
\end{table}

\vspace{0.2cm}

\begin{table}[h!]
\centering
\begin{tabularx}{\textwidth}{|X|}
\hline
\cellcolor[HTML]{F2F2F2}\textbf{Propósito} \\ \hline
Permitir al responsable poder registrar al cliente sus accesos a las actividades. \\ \hline
\end{tabularx}
\end{table}

\vspace{0.2cm}

\begin{table}[h!]
\centering
\begin{tabularx}{\textwidth}{|X|}
\hline
\cellcolor[HTML]{F2F2F2}\textbf{Resumen} \\ \hline
El responsable registra el acceso de un cliente a cierta actividad. \\ \hline
\end{tabularx}
\end{table}

\vspace{0.2cm}

\begin{table}[h!]
\centering
\begin{tabularx}{\textwidth}{|c|X|X|}
\hline
\multicolumn{3}{|c|}{\cellcolor[HTML]{F2F2F2}\textbf{Curso Normal (Básico)}} \\ \hline
1 & El responsable selecciona la opción para registrar acceso a la actividad. & \\ \hline
2 & & El sistema muestra el conjunto de actividades a las que es posible hacer el acceso. \\ \hline
3 & El responsable elige la actividad del cliente. & \\ \hline
4 & & El sistema pide el ID del cliente. \\ \hline
5 & El responsable introduce las credenciales del cliente. & \\ \hline
6 & & El sistema registra el acceso del cliente a la actividad. \\ \hline
\end{tabularx}
\end{table}

\vspace{0.2cm}

\begin{table}[h!]
\centering
\begin{tabularx}{\textwidth}{|c|X|}
\hline
\multicolumn{2}{|c|}{\cellcolor[HTML]{F2F2F2}\textbf{Cursos Alternos}} \\ \hline
3a & Si la actividad que solicita el cliente no está disponible, se le deberá de informar al cliente. \\ \hline
5a & Si el cliente no está registrado previamente, deberá rellenar la información necesaria primero. \\ \hline
\end{tabularx}
\end{table}

\vspace{0.2cm}

\begin{table}[h!]
\centering
\begin{tabularx}{\textwidth}{|X|X|X|X|}
\hline
\multicolumn{4}{|c|}{\cellcolor[HTML]{F2F2F2}\textbf{Otros datos}} \\ \hline
\textbf{Frecuencia esperada} &  & \textbf{Rendimiento} &  \\ \hline
\textbf{Importancia} &  & \textbf{Urgencia} &  \\ \hline
\textbf{Estado} &  & \textbf{Estabilidad} &  \\ \hline
\end{tabularx}
\end{table}

\vspace{0.2cm}

\begin{table}[h!]
\centering
\begin{tabularx}{\textwidth}{|X|}
\hline
\cellcolor[HTML]{F2F2F2}\textbf{Comentarios} \\ \hline
\\ % Aquí no había comentarios en la imagen, lo dejamos vacío.
\hline
\end{tabularx}
\end{table}

%======================================== Caso de uso 30 ====================

\begin{table}[h!]
\centering
\begin{tabularx}{\textwidth}{|l|X|X|}
\hline
\textbf{Caso de Uso} & Registrar fin de actividad & CU-30 \\ \hline
\textbf{Actores} & \multicolumn{2}{l|}{Responsable (I)} \\ \hline
\textbf{Tipo} & \multicolumn{2}{l|}{Primario, Esencial} \\ \hline
\textbf{Referencias} & \multicolumn{2}{l|}{} \\ \hline
\textbf{Precondición} & \multicolumn{2}{l|}{La actividad debe estar corriendo en el sistema en ese momento.} \\ \hline
\textbf{Poscondición} & \multicolumn{2}{l|}{Se registrará el fin de dicha actividad.} \\ \hline
\textbf{Autor} & \multicolumn{2}{l|}{Julian Carrion Tovar \hfill Fecha \hfill Versión 1.0} \\ \hline
\end{tabularx}
\end{table}

\vspace{0.2cm}

\begin{table}[h!]
\centering
\begin{tabularx}{\textwidth}{|X|}
\hline
\cellcolor[HTML]{F2F2F2}\textbf{Propósito} \\ \hline
Permitir al responsable poner fin a la actividad. \\ \hline
\end{tabularx}
\end{table}

\vspace{0.2cm}

\begin{table}[h!]
\centering
\begin{tabularx}{\textwidth}{|X|}
\hline
\cellcolor[HTML]{F2F2F2}\textbf{Resumen} \\ \hline
El responsable podrá poner fin a la actividad y no dejar acceso a la misma después de esto. \\ \hline
\end{tabularx}
\end{table}

\vspace{0.2cm}

\begin{table}[h!]
\centering
\begin{tabularx}{\textwidth}{|c|X|X|}
\hline
\multicolumn{3}{|c|}{\cellcolor[HTML]{F2F2F2}\textbf{Curso Normal (Básico)}} \\ \hline
1 & El responsable selecciona la opción de registrar el fin de la actividad. & \\ \hline
2 & & El sistema muestra las actividades que se encuentran en estado corriendo. \\ \hline
3 & El responsable selecciona la actividad que va a poner fin. & \\ \hline
4 & & El sistema registra el fin de la actividad, cambia su estado a finalizado y notifica al cliente. \\ \hline
\end{tabularx}
\end{table}

\vspace{0.2cm}

\begin{table}[h!]
\centering
\begin{tabularx}{\textwidth}{|c|X|}
\hline
\multicolumn{2}{|c|}{\cellcolor[HTML]{F2F2F2}\textbf{Cursos Alternos}} \\ \hline
1a & Si el responsable descubre que la actividad no se encuentra en curso en ese momento, el proceso se detendrá. \\ \hline
\end{tabularx}
\end{table}

\vspace{0.2cm}

\begin{table}[h!]
\centering
\begin{tabularx}{\textwidth}{|X|X|X|X|}
\hline
\multicolumn{4}{|c|}{\cellcolor[HTML]{F2F2F2}\textbf{Otros datos}} \\ \hline
\textbf{Frecuencia esperada} &  & \textbf{Rendimiento} &  \\ \hline
\textbf{Importancia} &  & \textbf{Urgencia} &  \\ \hline
\textbf{Estado} & & \textbf{Estabilidad} &  \\ \hline
\end{tabularx}
\end{table}

\vspace{0.2cm}

\begin{table}[h!]
\centering
\begin{tabularx}{\textwidth}{|X|}
\hline
\cellcolor[HTML]{F2F2F2}\textbf{Comentarios} \\ \hline
\\ % En la imagen no había comentarios, lo dejamos en blanco
\hline
\end{tabularx}
\end{table}


%======================================== Caso de uso 31 ====================


\clearpage

\begin{table}[h!]
\centering
\begin{tabularx}{\textwidth}{|p{3cm}|p{8cm}|p{2cm}|}
\hline
\textbf{Caso de Uso} & Informar de Incidencia & CU-31 \\ \hline
\textbf{Actores} & \multicolumn{2}{p{12cm}|}{Responsable} \\ \hline
\textbf{Tipo} & \multicolumn{2}{p{12cm}|}{Primario, Esencial} \\ \hline
\textbf{Referencias} & \multicolumn{2}{p{12cm}|}{} \\ \hline
\textbf{Precondición} & \multicolumn{2}{p{12cm}|}{La actividad debe estar registrada en el sistema y debe existir una incidencia real.} \\ \hline
\textbf{Poscondición} & \multicolumn{2}{p{12cm}|}{La incidencia queda registrada en el sistema para su posterior procesamiento.} \\ \hline

\textbf{Autor} & \multicolumn{2}{p{10cm}|}{Julian Carrion Tovar \hfill Fecha \hfill Versión 1.0} \\ \hline
\end{tabularx}
\end{table}


\vspace{0.2cm}

\begin{table}[h!]
\centering
\begin{tabularx}{\textwidth}{|X|}
\hline
\cellcolor[HTML]{F2F2F2}\textbf{Propósito} \\ \hline
Permitir al responsable notificar sobre la incidencia de la actividad. \\ \hline
\end{tabularx}
\end{table}

\vspace{0.2cm}

\begin{table}[h!]
\centering
\begin{tabularx}{\textwidth}{|X|}
\hline
\cellcolor[HTML]{F2F2F2}\textbf{Resumen} \\ \hline
El responsable puede informar sobre la incidencia que ha ocurrido en la realización de la actividad lo que permite al sistema gestionar la situación según los requerimientos. \\ \hline
\end{tabularx}
\end{table}

\vspace{0.2cm}

\begin{table}[h!]
\centering
\begin{tabularx}{\textwidth}{|c|X|X|}
\hline
\multicolumn{3}{|c|}{\cellcolor[HTML]{F2F2F2}\textbf{Curso Normal (Básico)}} \\ \hline
1 & El responsable selecciona la opción de informar de una incidencia en una actividad. & \\ \hline
2 & & El sistema presenta una variedad de opciones donde el responsable puede ir cumplimentándolas hasta describir la situación perfecta. \\ \hline
3 & El responsable completa la información de las diferentes opciones. & \\ \hline
4 & & El sistema registra toda la información necesaria y genera un ID para la incidencia. \\ \hline
\end{tabularx}
\end{table}

\vspace{0.2cm}

\begin{table}[h!]
\centering
\begin{tabularx}{\textwidth}{|c|X|}
\hline
\multicolumn{2}{|c|}{\cellcolor[HTML]{F2F2F2}\textbf{Cursos Alternos}} \\ \hline
3a & Si las opciones no son cumplimentadas correctamente, el sistema generará un mensaje de error. \\ \hline
\end{tabularx}
\end{table}

\vspace{0.2cm}

\begin{table}[h!]
\centering
\begin{tabularx}{\textwidth}{|X|X|X|X|}
\hline
\multicolumn{4}{|c|}{\cellcolor[HTML]{F2F2F2}\textbf{Otros datos}} \\ \hline
\textbf{Frecuencia esperada} &  & \textbf{Rendimiento} &  \\ \hline
\textbf{Importancia} &  & \textbf{Urgencia} &  \\ \hline
\textbf{Estado} & & \textbf{Estabilidad} &  \\ \hline
\end{tabularx}
\end{table}

\vspace{0.2cm}

\begin{table}[h!]
\centering
\begin{tabularx}{\textwidth}{|X|}
\hline
\cellcolor[HTML]{F2F2F2}\textbf{Comentarios} \\ \hline
Las incidencias deben de ser algo poco común si nuestro sistema funciona correctamente. \\ \hline
\end{tabularx}
\end{table}


\chapter{Gestión de Acceso de Usuarios}


%======================================== Caso de uso 32 ====================

\begin{table}[h!]
\centering
\begin{tabularx}{\textwidth}{|p{3cm}|p{8cm}|p{2cm}|}
\hline
\textbf{Caso de Uso} & Login & CU-32 \\ \hline
\textbf{Actores} & \multicolumn{2}{p{12cm}|}{Usuario} \\ \hline
\textbf{Tipo} & \multicolumn{2}{p{12cm}|}{Primario, Esencial} \\ \hline
\textbf{Referencias} & \multicolumn{2}{p{12cm}|}{} \\ \hline
\textbf{Precondición} & \multicolumn{2}{p{12cm}|}{Para hacer login, el usuario debe de estar registrado previamente en el sistema.} \\ \hline
\textbf{Poscondición} & \multicolumn{2}{p{12cm}|}{El usuario accede al sistema correctamente y se le informa.} \\ \hline
\textbf{Autor} & \multicolumn{2}{p{10cm}|}{Julian Carrión Tovar \hfill Fecha \hfill Versión 1.0} \\ \hline
\end{tabularx}
\end{table}

\vspace{0.2cm}

\begin{table}[h!]
\centering
\begin{tabularx}{\textwidth}{|X|}
\hline
\cellcolor[HTML]{F2F2F2}\textbf{Propósito} \\ \hline
Permitir que el usuario pueda acceder al sistema a través de sus datos. \\ \hline
\end{tabularx}
\end{table}

\vspace{0.2cm}

\begin{table}[h!]
\centering
\begin{tabularx}{\textwidth}{|X|}
\hline
\cellcolor[HTML]{F2F2F2}\textbf{Resumen} \\ \hline
El usuario proporciona sus datos y contraseña para realizar la operación de login y acceder al sistema. \\ \hline
\end{tabularx}
\end{table}

\vspace{0.2cm}

\begin{table}[h!]
\centering
\begin{tabularx}{\textwidth}{|c|X|X|}
\hline
\multicolumn{3}{|c|}{\cellcolor[HTML]{F2F2F2}\textbf{Curso Normal (Básico)}} \\ \hline
1 & El usuario accede a la página de login del sistema. & \\ \hline
2 & & El sistema presenta una página de inicio, donde puede ingresar los datos pertinentes. \\ \hline
3 & El usuario introduce sus datos y contraseña. & \\ \hline
4 & & El sistema verifica los datos y autentica al usuario. \\ \hline
5 & & El sistema asigna al usuario su cuenta. \\ \hline
\end{tabularx}
\end{table}

\vspace{0.2cm}

\begin{table}[h!]
\centering
\begin{tabularx}{\textwidth}{|c|X|}
\hline
\multicolumn{2}{|c|}{\cellcolor[HTML]{F2F2F2}\textbf{Cursos Alternos}} \\ \hline
4a & Si en la verificación los datos son incorrectos, el sistema deberá de generar un mensaje de error. \\ \hline
\end{tabularx}
\end{table}

\vspace{0.2cm}

\begin{table}[h!]
\centering
\begin{tabularx}{\textwidth}{|X|X|X|X|}
\hline
\multicolumn{4}{|c|}{\cellcolor[HTML]{F2F2F2}\textbf{Otros datos}} \\ \hline
\textbf{Frecuencia esperada} &  & \textbf{Rendimiento} &  \\ \hline
\textbf{Importancia} &  & \textbf{Urgencia} &  \\ \hline
\textbf{Estado} &  & \textbf{Estabilidad} &  \\ \hline
\end{tabularx}
\end{table}

\vspace{0.2cm}

\begin{table}[h!]
\centering
\begin{tabularx}{\textwidth}{|X|}
\hline
\cellcolor[HTML]{F2F2F2}\textbf{Comentarios} \\ \hline
La operación de login es esencial para que el sistema pueda permitir el acceso de los usuarios, y debe de ser eficaz y rápida. \\ \hline
\end{tabularx}
\end{table}

%======================================== Caso de uso 33 ====================

\clearpage

\begin{table}[h!]
\centering
\begin{tabularx}{\textwidth}{|p{3cm}|p{8cm}|p{2cm}|}
\hline
\textbf{Caso de Uso} & Logout & CU-33 \\ \hline
\textbf{Actores} & \multicolumn{2}{p{12cm}|}{Usuario} \\ \hline
\textbf{Tipo} & \multicolumn{2}{p{12cm}|}{Primario, Esencial} \\ \hline
\textbf{Referencias} & \multicolumn{2}{p{12cm}|}{} \\ \hline
\textbf{Precondición} & \multicolumn{2}{p{12cm}|}{El usuario debe de haber previamente iniciado sesión (login).} \\ \hline
\textbf{Poscondición} & \multicolumn{2}{p{12cm}|}{Al cerrar sesión el usuario debe de ser redirigido a la página de login.} \\ \hline
\textbf{Autor} & \multicolumn{2}{p{10cm}|}{Julian Carrión Tovar \hfill Fecha \hfill Versión 1.0} \\ \hline
\end{tabularx}
\end{table}

\vspace{0.2cm}

\begin{table}[h!]
\centering
\begin{tabularx}{\textwidth}{|X|}
\hline
\cellcolor[HTML]{F2F2F2}\textbf{Propósito} \\ \hline
Permitir que un usuario realice la operación de logout en el sistema (cierre de sesión). \\ \hline
\end{tabularx}
\end{table}

\vspace{0.2cm}

\begin{table}[h!]
\centering
\begin{tabularx}{\textwidth}{|X|}
\hline
\cellcolor[HTML]{F2F2F2}\textbf{Resumen} \\ \hline
El usuario puede cerrar sesión lo que hace que tenga que ingresar los datos necesarios para login de nuevo, si quiere volver a entrar, esta operación le permite la salida del sistema. \\ \hline
\end{tabularx}
\end{table}

\vspace{0.2cm}

\begin{table}[h!]
\centering
\begin{tabularx}{\textwidth}{|c|X|X|}
\hline
\multicolumn{3}{|c|}{\cellcolor[HTML]{F2F2F2}\textbf{Curso Normal (Básico)}} \\ \hline
1 & El usuario selecciona la opción de logout en el sistema. & \\ \hline
2 & & El sistema cierra la sesión del usuario. \\ \hline
3 & & El sistema lo redirige a la página de login. \\ \hline
\end{tabularx}
\end{table}

\vspace{0.2cm}

\begin{table}[h!]
\centering
\begin{tabularx}{\textwidth}{|c|X|}
\hline
\multicolumn{2}{|c|}{\cellcolor[HTML]{F2F2F2}\textbf{Cursos Alternos}} \\ \hline
 & \\ \hline
\end{tabularx}
\end{table}

\vspace{0.2cm}

\begin{table}[h!]
\centering
\begin{tabularx}{\textwidth}{|X|X|X|X|}
\hline
\multicolumn{4}{|c|}{\cellcolor[HTML]{F2F2F2}\textbf{Otros datos}} \\ \hline
\textbf{Frecuencia esperada} &  & \textbf{Rendimiento} &  \\ \hline
\textbf{Importancia} &  & \textbf{Urgencia} &  \\ \hline
\textbf{Estado} &  & \textbf{Estabilidad} &  \\ \hline
\end{tabularx}
\end{table}

\vspace{0.2cm}

\begin{table}[h!]
\centering
\begin{tabularx}{\textwidth}{|X|}
\hline
\cellcolor[HTML]{F2F2F2}\textbf{Comentarios} \\ \hline
Suponemos que un usuario cada vez que hace un login normalmente, tendrá que hacer una operación de logout por lo tanto la frecuencia será similar o equivalente. \\ \hline
\end{tabularx}
\end{table}




%======================================== Caso de uso 34 ====================
\clearpage
\chapter{Gestión de Responsables de Actividades}

\begin{table}[h!]
\centering
\begin{tabularx}{\textwidth}{|l|X|X|}
\hline
\textbf{Caso de Uso} & Alta responsable & CU-34 \\ \hline
\textbf{Actores} & \multicolumn{2}{l|}{Empleado(I), Responsable} \\ \hline
\textbf{Tipo} & \multicolumn{2}{l|}{Primario, Esencial} \\ \hline
\textbf{Referencias} & \multicolumn{2}{l|}{} \\ \hline
\textbf{Precondición} & \multicolumn{2}{l|}{El empleado debe de haber iniciado sesión (login) en el sistema.} \\ \hline
\textbf{Poscondición} & \multicolumn{2}{l|}{El nuevo responsable queda registrado en el sistema.} \\ \hline
\textbf{Autor} & \multicolumn{2}{l|}{Julian Carrion Tovar \hfill Fecha \hfill Versión 1.0} \\ \hline
\end{tabularx}
\end{table}

\vspace{0.2cm}

\begin{table}[h!]
\centering
\begin{tabularx}{\textwidth}{|X|}
\hline
\cellcolor[HTML]{F2F2F2}\textbf{Propósito} \\ \hline
Permitir que un empleado pueda registrar a un nuevo responsable. \\ \hline
\end{tabularx}
\end{table}

\vspace{0.2cm}

\begin{table}[h!]
\centering
\begin{tabularx}{\textwidth}{|X|}
\hline
\cellcolor[HTML]{F2F2F2}\textbf{Resumen} \\ \hline
Un empleado accederá al sistema con sus credenciales y podrá agregar al sistema un nuevo responsable con los procedimientos pertinentes, comunicándole su alta y sus credenciales por correo electrónico. \\ \hline
\end{tabularx}
\end{table}

\vspace{0.2cm}

\begin{table}[h!]
\centering
\begin{tabularx}{\textwidth}{|c|X|X|}
\hline
\multicolumn{3}{|c|}{\cellcolor[HTML]{F2F2F2}\textbf{Curso Normal (Básico)}} \\ \hline
1 & El empleado inicia sesión y selecciona la opción de agregar responsable. & \\ \hline
2 & & El sistema da los permisos al empleado necesarios para agregar al responsable. \\ \hline
3 & & Le presenta los campos que debe de rellenar el empleado para poder agregar al responsable. \\ \hline
4 & El empleado rellena los datos pertinentes que le solicita el sistema para registrar al nuevo responsable. & \\ \hline
5 & & El sistema verifica los datos que han sido introducidos por el empleado y registra al responsable. \\ \hline
\end{tabularx}
\end{table}

\vspace{0.2cm}

\begin{table}[h!]
\centering
\begin{tabularx}{\textwidth}{|c|X|}
\hline
\multicolumn{2}{|c|}{\cellcolor[HTML]{F2F2F2}\textbf{Cursos Alternos}} \\ \hline
5a & Si el empleado ha introducido los datos incorrectamente, el sistema le dará un mensaje de error y retrocederá. \\ \hline
\end{tabularx}
\end{table}

\vspace{0.2cm}

\begin{table}[h!]
\centering
\begin{tabularx}{\textwidth}{|X|X|X|X|}
\hline
\multicolumn{4}{|c|}{\cellcolor[HTML]{F2F2F2}\textbf{Otros datos}} \\ \hline
\textbf{Frecuencia esperada} &  & \textbf{Rendimiento} &  \\ \hline
\textbf{Importancia} &  & \textbf{Urgencia} &  \\ \hline
\textbf{Estado} &  & \textbf{Estabilidad} &  \\ \hline
\end{tabularx}
\end{table}

\vspace{0.2cm}

\begin{table}[h!]
\centering
\begin{tabularx}{\textwidth}{|X|}
\hline
\cellcolor[HTML]{F2F2F2}\textbf{Comentarios} \\ \hline
\\ \hline
\end{tabularx}
\end{table}


%======================================== Caso de uso 35 ====================

\clearpage


\begin{table}[h!]
\centering
\begin{tabularx}{\textwidth}{|l|X|X|}
\hline
\textbf{Caso de Uso} & Baja responsable & CU-35 \\ \hline
\textbf{Actores} & \multicolumn{2}{l|}{Empleado(I), Responsable} \\ \hline
\textbf{Tipo} & \multicolumn{2}{l|}{Primario, Esencial} \\ \hline
\textbf{Referencias} & \multicolumn{2}{l|}{} \\ \hline
\textbf{Precondición} & \multicolumn{2}{l|}{El responsable debe de haber estado de alta previamente.} \\ \hline
\textbf{Poscondición} & \multicolumn{2}{l|}{El responsable es eliminado/dado de baja del sistema.} \\ \hline
\textbf{Autor} & \multicolumn{2}{l|}{Julian Carrion Tovar \hfill Fecha \hfill Versión 1.0} \\ \hline
\end{tabularx}
\end{table}

\vspace{0.2cm}

\begin{table}[h!]
\centering
\begin{tabularx}{\textwidth}{|X|}
\hline
\cellcolor[HTML]{F2F2F2}\textbf{Propósito} \\ \hline
Permitir que un empleado pueda realizar la baja de un responsable. \\ \hline
\end{tabularx}
\end{table}

\vspace{0.2cm}

\begin{table}[h!]
\centering
\begin{tabularx}{\textwidth}{|X|}
\hline
\cellcolor[HTML]{F2F2F2}\textbf{Resumen} \\ \hline
Un empleado podrá acceder a la opción de baja, dar de baja al responsable y consecutivamente eliminar todos los datos asociados al mismo. \\ \hline
\end{tabularx}
\end{table}

\vspace{0.2cm}

\begin{table}[h!]
\centering
\begin{tabularx}{\textwidth}{|c|X|X|}
\hline
\multicolumn{3}{|c|}{\cellcolor[HTML]{F2F2F2}\textbf{Curso Normal (Básico)}} \\ \hline
1 & El empleado selecciona la opción de dar de baja al responsable. & \\ \hline
2 & & El sistema presenta la lista de responsables que se encuentran activos. \\ \hline
3 & El empleado selecciona al responsable que desea dar de baja. & \\ \hline
4 & & El sistema muestra la confirmación de dar de baja al responsable. \\ \hline
5 & El empleado confirma la baja. & \\ \hline
6 & & El sistema da de baja al responsable. \\ \hline
\end{tabularx}
\end{table}

\vspace{0.2cm}

\begin{table}[h!]
\centering
\begin{tabularx}{\textwidth}{|c|X|}
\hline
\multicolumn{2}{|c|}{\cellcolor[HTML]{F2F2F2}\textbf{Cursos Alternos}} \\ \hline
3a & Si el empleado que desea dar de baja, no se encuentra registrado o de alta, recibirá un mensaje de error de que no puede darlo de baja. \\ \hline
\end{tabularx}
\end{table}

\vspace{0.2cm}

\begin{table}[h!]
\centering
\begin{tabularx}{\textwidth}{|X|X|X|X|}
\hline
\multicolumn{4}{|c|}{\cellcolor[HTML]{F2F2F2}\textbf{Otros datos}} \\ \hline
\textbf{Frecuencia esperada} &  & \textbf{Rendimiento} &  \\ \hline
\textbf{Importancia} &  & \textbf{Urgencia} &  \\ \hline
\textbf{Estado} &  & \textbf{Estabilidad} &  \\ \hline
\end{tabularx}
\end{table}

\vspace{0.2cm}

\begin{table}[h!]
\centering
\begin{tabularx}{\textwidth}{|X|}
\hline
\cellcolor[HTML]{F2F2F2}\textbf{Comentarios} \\ \hline
La confirmación de dar de baja al responsable, es por meros motivos de seguridad. \\ \hline
\end{tabularx}
\end{table}

%======================================== Caso de uso 36 ====================


\clearpage
\begin{table}[h!]
\centering
\begin{tabularx}{\textwidth}{|l|X|X|}
\hline
\textbf{Caso de Uso} & Asignar actividad & CU-36 \\ \hline
\textbf{Actores} & \multicolumn{2}{l|}{Empleado(I), Responsable} \\ \hline
\textbf{Tipo} & \multicolumn{2}{l|}{Primario, Esencial} \\ \hline
\textbf{Referencias} & \multicolumn{2}{l|}{} \\ \hline
\textbf{Precondición} & \multicolumn{2}{l|}{Deben haber actividades y responsables existentes.} \\ \hline
\textbf{Poscondición} & \multicolumn{2}{l|}{La actividad queda asignada a un responsable en el sistema.} \\ \hline
\textbf{Autor} & \multicolumn{2}{l|}{Julian Carrion Tovar \hfill Fecha \hfill Versión 1.0} \\ \hline
\end{tabularx}
\end{table}

\vspace{0.2cm}

\begin{table}[h!]
\centering
\begin{tabularx}{\textwidth}{|X|}
\hline
\cellcolor[HTML]{F2F2F2}\textbf{Propósito} \\ \hline
Permitir que un empleado pueda asignar una actividad a un responsable específico. \\ \hline
\end{tabularx}
\end{table}

\vspace{0.2cm}

\begin{table}[h!]
\centering
\begin{tabularx}{\textwidth}{|X|}
\hline
\cellcolor[HTML]{F2F2F2}\textbf{Resumen} \\ \hline
El empleado selecciona la actividad y el responsable que quiere asignarle, y los asocia por medio del sistema. \\ \hline
\end{tabularx}
\end{table}

\vspace{0.2cm}

\begin{table}[h!]
\centering
\begin{tabularx}{\textwidth}{|c|X|X|}
\hline
\multicolumn{3}{|c|}{\cellcolor[HTML]{F2F2F2}\textbf{Curso Normal (Básico)}} \\ \hline
1 & El empleado selecciona la opción de asignar actividad a responsable. & \\ \hline
2 & & El sistema le muestra las actividades y responsables disponibles en ese momento. \\ \hline
3 & El empleado selecciona la actividad. & \\ \hline
4 & El empleado selecciona el responsable. & \\ \hline
5 & & El sistema asocia al responsable a dicha actividad. \\ \hline
\end{tabularx}
\end{table}

\vspace{0.2cm}

\begin{table}[h!]
\centering
\begin{tabularx}{\textwidth}{|c|X|}
\hline
\multicolumn{2}{|c|}{\cellcolor[HTML]{F2F2F2}\textbf{Cursos Alternos}} \\ \hline
3a & Si no hubiera actividades disponibles, se mandaría un mensaje de error. \\ \hline
4a & Si no hubiera responsables disponibles o adecuados, se mandaría un mensaje de error, o se volvería a la selección de un nuevo responsable respectivamente. \\ \hline
\end{tabularx}
\end{table}

\vspace{0.2cm}

\begin{table}[h!]
\centering
\begin{tabularx}{\textwidth}{|X|X|X|X|}
\hline
\multicolumn{4}{|c|}{\cellcolor[HTML]{F2F2F2}\textbf{Otros datos}} \\ \hline
\textbf{Frecuencia esperada} &  & \textbf{Rendimiento} &  \\ \hline
\textbf{Importancia} &  & \textbf{Urgencia} &  \\ \hline
\textbf{Estado} &  & \textbf{Estabilidad} &  \\ \hline
\end{tabularx}
\end{table}

\vspace{0.2cm}

\begin{table}[h!]
\centering
\begin{tabularx}{\textwidth}{|X|}
\hline
\cellcolor[HTML]{F2F2F2}\textbf{Comentarios} \\ \hline
\\ \hline
\end{tabularx}
\end{table}



%======================================== Caso de uso 37 ====================
\clearpage


\begin{table}[h!]
\centering
\begin{tabularx}{\textwidth}{|p{3cm}|p{8cm}|p{2cm}|}
\hline
\textbf{Caso de Uso} & Cambiar responsable de actividad & CU-37 \\ \hline
\textbf{Actores} & \multicolumn{2}{p{12cm}|}{Empleado(I), Responsable} \\ \hline
\textbf{Tipo} & \multicolumn{2}{p{12cm}|}{Primario, Esencial} \\ \hline
\textbf{Referencias} & \multicolumn{2}{p{12cm}|}{} \\ \hline
\textbf{Precondición} & \multicolumn{2}{p{12cm}|}{La actividad debe de tener un responsable y debe haber otros disponibles.} \\ \hline
\textbf{Poscondición} & \multicolumn{2}{p{12cm}|}{El nuevo responsable quedará asignado a su nueva actividad.} \\ \hline
\textbf{Autor} & \multicolumn{2}{p{10cm}|}{Julian Carrion Tovar \hfill Fecha \hfill Versión 1.0} \\ \hline
\end{tabularx}
\end{table}

\vspace{0.2cm}

\begin{table}[h!]
\centering
\begin{tabularx}{\textwidth}{|X|}
\hline
\cellcolor[HTML]{F2F2F2}\textbf{Propósito} \\ \hline
Permitir que un empleado pueda cambiar el responsable que se ocupa de cierta actividad. \\ \hline
\end{tabularx}
\end{table}

\vspace{0.2cm}

\begin{table}[h!]
\centering
\begin{tabularx}{\textwidth}{|X|}
\hline
\cellcolor[HTML]{F2F2F2}\textbf{Resumen} \\ \hline
Un empleado selecciona la opción de cambiar responsable, y realiza los cambios pertinentes en la organización del sistema para que un nuevo responsable se ocupe de una nueva actividad. Este cambio es comunicado al responsable anterior y al nuevo responsable mediante correo electrónico. Al nuevo responsable se le indica también los datos de la actividad. \\ \hline
\end{tabularx}
\end{table}

\vspace{0.2cm}

\begin{table}[h!]
\centering
\begin{tabularx}{\textwidth}{|c|X|X|}
\hline
\multicolumn{3}{|c|}{\cellcolor[HTML]{F2F2F2}\textbf{Curso Normal (Básico)}} \\ \hline
1 & El empleado selecciona la opción de cambiar responsable de actividad. & \\ \hline
2 & & El sistema muestra la lista de actividades y responsables. \\ \hline
3 & El empleado selecciona la actividad que quiere cambiar el responsable. & \\ \hline
4 & & El sistema le muestra los responsables disponibles. \\ \hline
5 & El empleado cambia el responsable de dicha actividad. & \\ \hline
6 & & El sistema realiza el cambio de responsable de actividad y actualiza la información, posteriormente lo notifica. \\ \hline
\end{tabularx}
\end{table}

\vspace{0.2cm}

\begin{table}[h!]
\centering
\begin{tabularx}{\textwidth}{|c|X|}
\hline
\multicolumn{2}{|c|}{\cellcolor[HTML]{F2F2F2}\textbf{Cursos Alternos}} \\ \hline
4a & Si no hay responsables disponibles, el sistema muestra un mensaje de error y retrocede al paso 1 posteriormente. \\ \hline
\end{tabularx}
\end{table}

\vspace{0.2cm}

\begin{table}[h!]
\centering
\begin{tabularx}{\textwidth}{|X|X|X|X|}
\hline
\multicolumn{4}{|c|}{\cellcolor[HTML]{F2F2F2}\textbf{Otros datos}} \\ \hline
\textbf{Frecuencia esperada} &  & \textbf{Rendimiento} &  \\ \hline
\textbf{Importancia} &  & \textbf{Urgencia} &  \\ \hline
\textbf{Estado} &  & \textbf{Estabilidad} &  \\ \hline
\end{tabularx}
\end{table}

\vspace{0.2cm}

\begin{table}[h!]
\centering
\begin{tabularx}{\textwidth}{|X|}
\hline
\cellcolor[HTML]{F2F2F2}\textbf{Comentarios} \\ \hline
El cambio de responsable en una actividad, considero que es una situación extraordinaria cuando existe algún tipo de incidencia. \\ \hline
\end{tabularx}
\end{table}


